\head{两种载体的同文校订与真实接续点}
\source{\texttt{EGAV1-2.ps}完整六页；\texttt{EGAVaiello.pdf}第62--66页；Dot随后两批第16节译稿}

\audit{D17-S-C01}{重合范围与未重复计数}{两批作为校订入编，保留既有全文，接续点重新据原件核定}
第十六版的\texttt{text/v1\_2.tex}已经连续翻译了PS载体的全部可读文字，包括定义1、备注2、定义3、备注4(a)--(d)、命题5、二次初始项及判别式构造、命题7和一般局部自由模情形的讨论。
Dot在总包之后交付的两批第16节译稿，只覆盖同一短篇的定义1至备注4(c)。本版将其相应校订纳入旧短篇，不重复追加一份正文；原有备注4(d)及随后全部文字均保留。

两种载体的末端不同。PS在“有下面这个图的交换性”一句后结束；合并英文稿第66页在同一句之后确实继续给出交换图、符号说明和论证。故第62页备注4(d)只是Dot分支当时的断点，\textbf{本工程当前第V卷的真正缺译接续点是第66页该交换图}。第IV章主体的接续点仍为\texttt{EGAIV.1.pdf}第221页的内容提要，随后章首和第223页的（1.1.1）。

本版未将合并英文稿拟定目录里的第16节“转换为线性系的语言”与第62页Part II的“§16 New EGA V”混为一节；前者是第一篇的原拟定计划，后者是另篇预备笔记的现有载体标题。

\audit{D17-S-C02}{旧短篇已核定错误移入正确正文}{类型、基底、维数与Jacobian判据的局部修订}
Dot两批提供的主要订正见本报告V-H-LATE-06。除此以外，本次把旧短篇中已经列在T01、T08--T10及E06等脚注中的若干确定问题落实到正文；原错式、原措辞和改动依据均留在D17-S01--S13脚注。未改写第I、II、III章正文。

\begin{longtable}{@{}p{.16\linewidth}p{.76\linewidth}@{}}
\toprule 注号 & 正文处理\\\midrule\endhead
S01--S02 & 命题5及二次初始项的基变换，使用本段已经给定的相对底$Y$，原文误写$S$。\\
S03 & 相对维数计算为$(m+d)-(d+1)=m-1$；原式$(m+1)-(d+1)=m-1$不成立。\\
S04--S06 & 零点概形的微分列首项限制到$W$；正确Jacobian条件是余域纤维单射，且光滑相对底为$S$；余法正合列末项为$\Omega^1_{W/S}$。\\
S07--S08 & 图中主部及切层的相对下标为$X/Y$；将原错置的$M,V$恢复为图示$\mu,\nu$，并区分两个商模的满射性与二次型的特征问题。\\
S09 & 外幂取目标秩$r$：$\bigwedge^rQ\to\bigwedge^rM$；对应截面属于$(\bigwedge^rQ)^\vee\otimes\det M$。原根号状破损公式保留在脚注。\\
S10--S11 & 一般线性代数情形用$p+q=m+1$、$M/Q$及$L=\det P\otimes\det Q\otimes(\det M)^{-1}$，原错式逐项留注。\\
S12 & 把后续原文实际使用的局部直和因子条件提前列在使用商模行列式之前；单有局部自由子模不足以保证商模局部自由。\\
S13 & 根据作者在合并英文稿第73--74页的明确修订，把旧称“普通奇异零点”校为“奇异而非超奇异的零点”；不把特征2中的普通二次奇点等同于非超奇异零点。\\
\bottomrule
\end{longtable}

这些处理有直接的类型检验。对秩$r$的目标模$M$，满射性可在局部基中以$r$阶子式检验，因而对应$\bigwedge^rQ\to\det M$。在一般情形$p+q=m+1$下，$S=M/Q$秩为$p-1$、$R=M/P$秩为$q-1$，所以
\[
 (\bigwedge^{p-1}P)^\vee\otimes\det S
 \simeq P\otimes(\det P)^{-1}\otimes\det M\otimes(\det Q)^{-1},
\]
正是$\mathscr{H}\!om(L,P)$；$Q$侧同理。这说明$L$中$\det M$必须取逆，并且各商模须为局部自由模。

对Jacobian处的订正，局部写$F\simeq\mathcal O_X^c$，截面由$c$个函数给出。其微分在余域纤维上线性无关，等价于相应微分矩阵有可逆的$c$阶子式；在$X/S$光滑的环境中，这正是零点概形对$S$光滑且具有相对余维$c$的判据。它比模层同态抽象单射强；旧报告E06已经用尖点区别这两者，本版保留该历史反例，不另声称所有相关草稿论证已被修复。

\audit{D17-S-C03}{局部非零条件与特征2术语}{复核后在正文落实，并记录后续原文出处}
注2的末句必须按局部概形解释。若把“在某个邻域内$V(\varphi)\ne X$”只读作普通存在量词，取
\[
 X=\mathbb A^1_k\amalg\mathbb A^1_k,\qquad \varphi=(0,1),
\]
并取$x$在第一分支，则$V(\varphi)$在$x$处正则，且在邻域$X$内不等于$X$，但$\varphi_x=0$。因此该邻域措辞不足以刻画非奇异零点。本版正文改为在$x$处的局部子概形不等；等价地，局部环中的芽$\varphi_x$非零。此处区分的是局部条件与任意较大邻域上的不等，并不是更改原文已经写明的芽非零条件。

合并英文稿第73--74页有原作者对“普通奇异零点”旧称的明确修订。其特征2例子$X^2+YZ=0$说明，通常的普通二次奇点与极化双线性型非退化不是同一条件：三变量交替双线性型退化，但该锥顶仍是通常意义的普通二次奇点。第78页（16.1）--（16.2）明确给出相关对称双线性型及判别式$d(Q)$，并区分“非退化二次型”与“光滑二次型”。本版已经把这一明确的术语修订用于定义1末尾，并在脚注保留旧称及出处。第73--74页和第78页在本轮只为定点核查而读取，没有计作连续翻译完成。

\audit{D17-S-P01}{仍未解决的草稿问题}{同文校订不等于完整数学审定}
特征2的术语问题已有上列后续原文依据；本版仍未提前改写整个后续定义体系，也没有把行列式计算或作者对旧稿形式正确性的评价当作全部证明的独立审定。原稿尚未正式列出的“命题6”仍不补造，命题7所引用的“以上条件”及第66页起的概形结构论证仍须随下一段继续核查。

旧短篇原源码逐字保存在\path{history/v16_short_before_integration.tex}；本次两批Dot源码、各自进度及原始工程包均在交接包的原始交付档案中。历史报告的旧注号继续作为历史索引，不据新排版推定所有旧意见已经全部复核。
